% \iffalse meta-comment
%
% banshi-preset.dtx 是 Word 新建空白文档的那一套默认值
%
% \fi
%
% \section{预设：Word 的空白文档}
%
% \changes{v0.1.0}{2026/09/23}{加包选项 \option{preset}，默认 \texttt{blank}：
%   单独加载本包即照 Word 新建空白文档排}
%
% 在 \cls{ctexart} 或 \cls{article} 上写一句 \cs{usepackage}\marg{banshi}，排出来就是
% Word 里“新建空白文档”的样子：页面、字体、字号、行距、段后、对齐都照它的
% \file{Normal.dotm}。文档类自己管版面、只想借本包的机制时写
% \cs{usepackage}\oarg{preset=none}\marg{banshi}，本节什么都不做。
%
% 取值抄自 Word 16（2024 版）简体中文界面的 \file{Normal.dotm}，主题字体语言
% \texttt{zh-CN}：
% \begin{itemize}
% \item 纸张 A4；页边距上下 1440 缇（2.54\,cm）、左右 1800 缇（3.17\,cm）；页眉 851 缇、
%   页脚 992 缇；文档网格只定行（\texttt{type="lines"}），行距 312 缇即 15.6\,bp。
% \item 正文（\texttt{docDefaults} 加 Normal 样式）：等线，11\,bp，多倍行距
%   278/240（1.1583 倍），段后 160 缇即 8\,bp，左对齐，无首行缩进，孤行控制关。
% \item 标题 1 至 4：24、20、16、14\,bp，不加粗，颜色 \texttt{0F4761}；段前段后
%   24/4、8/4、8/4、4/2\,bp，行距随正文；前三级用主题的标题字体“等线 Light”，
%   第四级用正文字体。这四个样式在模板里是潜在的，要在 Word 里用过一次才写进
%   \file{styles.xml}，数是这样取出来的。
% \item 兼容模式 15。
% \end{itemize}
%
% 标题的层级对照 Word 的“标题 1”到“标题 4”：有 \cs{chapter} 的文档类从
% \cs{chapter} 算起，没有的从 \cs{section} 算起。标题要经 \pkg{ctex} 的
% \cs{ctexset} 发出去，\cls{article} 这类不带 \pkg{ctex} 的文档类只排正文，标题
% 照文档类原样。页码、页眉页脚也不动：Word 的空白文档没有页码，\LaTeX{} 的文档类
% 大多有，要去掉就写 \cs{pagestyle}\marg{empty}。
%
% 字体要 \hologo{LuaLaTeX} 或 \hologo{XeLaTeX}。系统里找不到等线就报一条警告，
% 字体照文档类原样，版面照样按等线的行高算。“等线 Light”在 macOS 上只随 Word
% 附带、不装进系统，找不到时前三级标题退回等线。
%
% ^^A ======================================================================
% ^^A Module banshi-preset: 包选项与空白文档预设
% ^^A ======================================================================
%    \begin{macrocode}
%<*banshi-preset>
%    \end{macrocode}
%
% 包选项只有这一个。其余设置写进 \cs{banshisetup}，放在 \cs{usepackage} 之后，
% 预设的取值就被它盖掉；做成包选项的话，预设在选项之后才装，反倒要盖掉用户写的。
%    \begin{macrocode}
\msg_new:nnnn { banshi } { preset-font-missing }
  { The~font~'#1'~was~not~found;~the~document~class~fonts~are~kept. }
  {
    The~blank~preset~uses~DengXian~(Word's~default~font~for~simplified~
    Chinese).~Line~heights~are~still~computed~from~DengXian's~metrics.
  }
\msg_new:nnnn { banshi } { preset-engine }
  { The~blank~preset~sets~fonts~only~under~LuaLaTeX~or~XeLaTeX. }
  { Page~setup~and~paragraph~spacing~still~apply;~fonts~are~left~alone. }
\tl_new:N \g__banshi_preset_tl
\tl_gset:Nn \g__banshi_preset_tl { blank }
\keys_define:nn { banshi / option }
  {
    preset .choices:nn = { blank , none }
      { \tl_gset_eq:NN \g__banshi_preset_tl \l_keys_choice_tl } ,
  }
%    \end{macrocode}
%
% \emph{页面。} 页眉那一段没有下框线，线宽与离字都写零。
%    \begin{macrocode}
\cs_new_protected:Npn \__banshi_preset_page:
  {
    \RequirePackage { geometry }
    \banshi_msword_layout:n
      {
        margin-top = 72bp , margin-bottom = 72bp ,
        margin-left = 90bp , margin-right = 90bp ,
        header = 42.55bp , footer = 49.6bp ,
        line-pitch = 15.6bp ,
        header-rule-width = 0bp , header-rule-from-text = 0bp ,
      }
  }
%    \end{macrocode}
%
% \emph{段落格式。} 目标名借 \texttt{chapter} 到 \texttt{subsubsection} 这四个，
% 分别存“标题 1”到“标题 4”，发的时候再按文档类有没有 \cs{chapter} 对到层级上。
% 行距 1.158333 是 278/240；网格只贴纵向，空白文档没有字符网格。
%    \begin{macrocode}
\cs_new_protected:Npn \__banshi_preset_format:
  {
    \banshi_format_set:n
      {
        body =
          {
            font-zh = dengxian , size = 11bp , line-spacing = 1.158333 ,
            space-after = 8bp , alignment = left ,
            snap-to-grid = { vertical = true , horizontal = false } ,
          } ,
        chapter =
          {
            font-zh = dengxian , size = 24bp , line-spacing = 1.158333 ,
            space-before = 24bp , space-after = 4bp , alignment = left ,
            snap-to-grid = { vertical = true , horizontal = false } ,
          } ,
        section =
          {
            font-zh = dengxian , size = 20bp , line-spacing = 1.158333 ,
            space-before = 8bp , space-after = 4bp , alignment = left ,
            snap-to-grid = { vertical = true , horizontal = false } ,
          } ,
        subsection =
          {
            font-zh = dengxian , size = 16bp , line-spacing = 1.158333 ,
            space-before = 8bp , space-after = 4bp , alignment = left ,
            snap-to-grid = { vertical = true , horizontal = false } ,
          } ,
        subsubsection =
          {
            font-zh = dengxian , size = 14bp , line-spacing = 1.158333 ,
            space-before = 4bp , space-after = 2bp , alignment = left ,
            snap-to-grid = { vertical = true , horizontal = false } ,
          } ,
      }
    \keys_set:nn { banshi } { word-compat = 15 }
    \clist_map_inline:nn { chapter , section , subsection }
      {
        \banshi_format_add_font:nn {##1}
          { \__banshi_preset_major: \__banshi_preset_color: }
      }
    \banshi_format_add_font:nn { subsubsection } { \__banshi_preset_color: }
  }
\cs_new_protected:Npn \__banshi_preset_color:
  { \cs_if_exist:NT \color_select:nn { \color_select:nn { HTML } { 0F4761 } } }
%    \end{macrocode}
%
% \emph{字体。} 正文的西文与汉字都是等线（Word 主题的正文字体两栏都填等线）。
% \cs{dengxian} 与 \pkg{ctex} 的 \cs{songti} 同类，只换汉字那一副，\option{font-zh}
% 写 \texttt{dengxian} 的目标发的就是它；找不到字体时它什么都不做。
%    \begin{macrocode}
\bool_new:N \g__banshi_preset_dengxian_bool
\cs_new_protected:Npn \__banshi_preset_major: { }
\cs_new_protected:Npn \__banshi_preset_fonts:
  {
    \bool_lazy_or:nnTF { \sys_if_engine_luatex_p: } { \sys_if_engine_xetex_p: }
      {
        \cs_if_exist:NF \setmainfont { \RequirePackage { fontspec } }
        \IfFontExistsTF { DengXian }
          {
            \bool_gset_true:N \g__banshi_preset_dengxian_bool
            \setmainfont { DengXian }
            \cs_if_exist:NT \setCJKmainfont { \setCJKmainfont { DengXian } }
            \cs_if_exist:NT \setCJKfamilyfont
              { \setCJKfamilyfont { dengxian } { DengXian } }
          }
          { \msg_warning:nnn { banshi } { preset-font-missing } { DengXian } }
        \IfFontExistsTF { DengXian~Light }
          {
            \newfontfamily \banshi@dengxian@light { DengXian~Light }
            \cs_if_exist:NTF \setCJKfamilyfont
              {
                \setCJKfamilyfont { dengxian-light } { DengXian~Light }
                \cs_gset_protected:Npn \__banshi_preset_major:
                  { \banshi@dengxian@light \CJKfamily { dengxian-light } }
              }
              { \cs_gset_protected:Npn \__banshi_preset_major: { \banshi@dengxian@light } }
          }
          { }
      }
      { \msg_warning:nn { banshi } { preset-engine } }
    \cs_if_exist:NF \dengxian
      {
        \cs_new_protected:Npn \dengxian
          {
            \bool_lazy_and:nnT
              { \g__banshi_preset_dengxian_bool }
              { \cs_if_exist_p:N \CJKfamily }
              { \CJKfamily { dengxian } }
          }
      }
  }
%    \end{macrocode}
%
% \emph{导言区结束时接上。} \cs{normalsize} 改成读正文的字号与行距，原来那一份
% 留着先执行，行间公式上下的间距仍是文档类给的。标题按有没有 \cs{chapter}
% 对层级：有就一一对应，没有就整体错开一格，“标题 1”发给 \texttt{section}。
%
% \cs{baselinestretch} 先设回 1：\pkg{ctex} 默认把它设成 1.3，行距是本包按 Word
% 算好的，再乘一次就成了 23.49\,bp（Word 是 18.07）。
%
% 汉字与西文紧挨着时中间的胶取 0.225 个字号：没有字符网格，Word 的“自动调整中文与
% 西文的间距”按字号给，11\,bp 下两侧量得 2.57 与 2.45。源码里汉字与西文之间敲了
% 空格的，Word 是空格再加这一截，\pkg{luatexja} 见了空格就不加，这里还差 2.5\,bp
% 左右，没补。
%
% 左对齐的段交回 \hologo{TeX} 断行，打开 \cs{banshi\_linebreak\_rigid\_ragged:}，
% 词隙与标点不缩，一行塞的字与 Word 一样多。段与段之间按行盒接，打开
% \cs{banshi\_format\_word\_interline:}。
%    \begin{macrocode}
\cs_new_protected:Npn \__banshi_preset_normalsize: { }
\cs_new_protected:Npn \__banshi_preset_apply:
  {
    \cs_set:Npn \baselinestretch { 1 }
    \cs_gset_eq:NN \__banshi_preset_normalsize: \normalsize
    \cs_gset_protected:Npn \normalsize
      {
        \__banshi_preset_normalsize:
        \@setfontsize \normalsize
          { \g_banshi_body_size_dim } { \g_banshi_body_lead_dim }
      }
    \banshi_format_word_interline:
    \banshi_linebreak_rigid_ragged:
    \banshi_format_apply_body:
    \banshi_format_hglue:nn { 0pt } { 0.225 \g_banshi_body_size_dim }
    \cs_if_exist:NT \ctexset
      {
        \cs_if_exist:NTF \chapter
          {
            \banshi_format_apply_headings:n
              { chapter , section , subsection , subsubsection }
          }
          {
            \banshi_format_apply_heading_to:nn { chapter }    { section }
            \banshi_format_apply_heading_to:nn { section }    { subsection }
            \banshi_format_apply_heading_to:nn { subsection } { subsubsection }
          }
      }
  }
\cs_new_protected:Npn \__banshi_preset_blank:
  {
    \__banshi_preset_page:
    \__banshi_preset_format:
    \__banshi_preset_fonts:
    \hook_gput_code:nnn { begindocument } { banshi } { \__banshi_preset_apply: }
  }
\ProcessKeyOptions [ banshi / option ]
\str_if_eq:VnT \g__banshi_preset_tl { blank } { \__banshi_preset_blank: }
%    \end{macrocode}
%
%    \begin{macrocode}
%</banshi-preset>
%    \end{macrocode}
%
% \Finale
